% Appendix F -- the staging/zero-copy crossover in full, moved out of section 8
% (2026-09-26) for the page limit. The body states the ordering and why each
% system is right for its own design point; this is the argument behind it.

\section{Staging versus zero-copy, in full}
\label{app:coherence}

Our host tier stages: a layer's KV is copied from a pinned host mirror into a
device rolling slot before attention reads it. DirectKV's kernel instead reads
CPU-resident KV directly, with no staging buffer. We built on a staged backend
for two practical reasons --- it runs on PCIe as well as on a coherent link,
which is what let us measure the negative regime at all, and their kernel is
SM90a-only with no build on our development hardware. Our results are therefore
\emph{conservative with respect to mechanism}: a better transfer path would
raise them.

It is still tempting to read the buffer as a tax that zero-copy removes, and
that reading is wrong in a way that decides which mechanism belongs where. The
two designs hide the transfer behind different things. A zero-copy kernel
overlaps its fetch with the arithmetic of the attention tile it is already
computing --- a window of one tile, and decode attention's intensity is the GQA
group size, 4--8 FLOP/byte against a balance point near 100, so there is almost
no arithmetic there to hide behind. Staging overlaps the fetch for layer
$\ell{+}1$ with the entire compute of layer $\ell$, including the QKV, output
and MLP projections, which are weight-bound and long --- 21 of a 34\,ms step for
Qwen2.5-32B by the projection of \S\ref{sec:model}. \textbf{The staging buffer
buys a layer-sized hiding window instead of a tile-sized one.}

What it costs is two extra HBM passes per byte and the slot's own footprint, so
the comparison is a crossover in offload ratio rather than a dominance. When the
per-step transfer fits inside the layer compute window, staging hides it and pays
only the extra HBM traffic, which at $\Xlink/\BW \approx 0.105$ is a small
fraction of something already hidden; when the transfer exceeds that window it is
exposed under either design, and zero-copy wins by avoiding the passes. This is
why the two systems make opposite choices and both are right: DirectKV places
100\% of KV on the host, so its per-step transfer is the whole working set and can
never fit a compute window, while we place a selected subset and size it to fit
by construction (\S\ref{sec:design:admission}). \textbf{The offload ratio is the
quantity our policy controls}, so the choice of mechanism follows from the policy
rather than constraining it.

A coherent link would help one cost that is not about hiding. Each step projects
new $K,V$ for every host-tier request and we copy them device-to-host to keep the
mirror authoritative; mapped host memory would let the kernel write them in
place, as DirectKV's Alg.~3 does. The saving is small in time --- an asynchronous
version measured 1.6\,ms --- but it removes a class of state rather than an
amount of time, and that state is where migration bugs live. What coherence does
not buy is a cheaper transport: it is a guarantee about visibility, and the same
bytes still cross the link at link bandwidth when the kernel touches them.

